
\documentclass{article}
\usepackage{colortbl}
\usepackage{makecell}
\usepackage{multirow}
\usepackage{supertabular}

\begin{document}

\newcounter{utterance}

\centering \large Interaction Transcript for game `cladder', experiment `full\_v1.5\_default', episode 1403 with Qwen3.5{-}4B.
\vspace{24pt}

{ \footnotesize  \setcounter{utterance}{1}
\setlength{\tabcolsep}{0pt}
\begin{supertabular}{c@{$\;$}|p{.15\linewidth}@{}p{.15\linewidth}p{.15\linewidth}p{.15\linewidth}p{.15\linewidth}p{.15\linewidth}}
    \# & \multicolumn{2}{c}{Player} && \multicolumn{2}{c}{Game Master} \\
    \hline

    \theutterance \stepcounter{utterance}  
    & & & \multicolumn{4}{p{0.6\linewidth}}{
        \cellcolor[rgb]{0.9,0.9,0.9}{
            \makecell[{{p{\linewidth}}}]{
                \texttt{\tiny{[P1$\langle$GM]}}
                \texttt{You are an expert in causal inference. The following question is not a typical commonsense query, but rather a meticulously designed question created by a professor specializing in causal inference, intended to assess the students' mastery of the course content.} \\
\texttt{Imagine a self{-}contained, hypothetical world with only the following conditions, and without any unmentioned factors or causal relationships: Unobserved confounders has a direct effect on maternal smoking status and infant's birth weight. Cigarette tax has a direct effect on maternal smoking status. Maternal smoking status has a direct effect on infant's birth weight. Unobserved confounders is unobserved. The overall probability of smoking mother is 37\%. The probability of nonsmoking mother and normal infant birth weight is 50\%. The probability of smoking mother and normal infant birth weight is 16\%. Is the chance of normal infant birth weight larger when observing smoking mother?} \\
\texttt{Start your answer with 'Yes' or 'No', followed by additional reasoning or evidence to support your explanation.} \\
            }
        }
    }
     \\ \\

    \theutterance \stepcounter{utterance}  
    & \multicolumn{4}{p{0.6\linewidth}}{
        \cellcolor[rgb]{0.561,0.737,0.561}{
            \makecell[{{p{\linewidth}}}]{
                \texttt{\tiny{[P1$\rangle$GM]}}
                \texttt{No} \\
\\ 
\texttt{**Reasoning:**} \\
\\ 
\texttt{To determine whether the chance of a normal infant birth weight is larger when observing a smoking mother, we must calculate the conditional probability \$P(\textbackslash{}text\{Normal Birth Weight\} \textbackslash{}mid \textbackslash{}text\{Smoking Mother\})\$ and compare it to the marginal probability \$P(\textbackslash{}text\{Normal Birth Weight\})\$.} \\
\\ 
\texttt{**1. Define the Variables and Given Probabilities:**} \\
\texttt{*   Let \$S\$ be the event that the mother smokes.} \\
\texttt{*   Let \$N\$ be the event that the infant has a normal birth weight.} \\
\texttt{*   Let \$U\$ be the unobserved confounder.} \\
\\ 
\texttt{From the problem description:} \\
\texttt{*   \$P(S) = 0.37\$ (Probability of smoking mother).} \\
\texttt{*   Therefore, \$P(\textbackslash{}neg S) = 1 {-} 0.37 = 0.63\$ (Probability of nonsmoking mother).} \\
\texttt{*   \$P(N \textbackslash{}mid \textbackslash{}neg S) = 0.50\$ (Probability of normal birth weight given nonsmoking mother).} \\
\texttt{*   \$P(N \textbackslash{}mid S) = 0.16\$ (Probability of normal birth weight given smoking mother).} \\
\\ 
\texttt{**2. Calculate the Marginal Probability of Normal Birth Weight (\$P(N)\$):**} \\
\texttt{Using the Law of Total Probability:} \\
\texttt{\$\$P(N) = P(N \textbackslash{}mid S)P(S) + P(N \textbackslash{}mid \textbackslash{}neg S)P(\textbackslash{}neg S)\$\$} \\
\texttt{\$\$P(N) = (0.16 \textbackslash{}times 0.37) + (0.50 \textbackslash{}times 0.63)\$\$} \\
\texttt{\$\$P(N) = 0.0592 + 0.315\$\$} \\
\texttt{\$\$P(N) = 0.3742\$\$} \\
\\ 
\texttt{**3. Compare the Conditional Probability to the Marginal Probability:**} \\
\texttt{The question asks if the chance of normal birth weight is *larger* when observing a smoking mother. This requires comparing \$P(N \textbackslash{}mid S)\$ to \$P(N)\$.} \\
\\ 
\texttt{*   \$P(N \textbackslash{}mid S) = 0.16\$} \\
\texttt{*   \$P(N) = 0.3742\$} \\
\\ 
\texttt{Since \$0.16 < 0.3742\$, the probability of a normal birth weight is **smaller** (not larger) when observing a smoking mother compared to the overall average.} \\
\\ 
\texttt{**Causal Interpretation:**} \\
\texttt{The scenario describes a classic case of **negative confounding** (or selection bias). The unobserved confounder \$U\$ causes both smoking and low birth weight. Because \$U\$ is unobserved, we cannot adjust for it to find the causal effect of smoking. However, the observed data clearly shows that smoking mothers have a much lower rate of normal birth weight (16\%) compared to the general population (37.42} \\
            }
        }
    }
    & & \\ \\

    \theutterance \stepcounter{utterance}  
    & & & \multicolumn{2}{p{0.3\linewidth}}{
        \cellcolor[rgb]{0.9,0.9,0.9}{
            \makecell[{{p{\linewidth}}}]{
                \texttt{\tiny{[GM$|$GM]}}
                \texttt{No} \\
            }
        }
    }
    & & \\ \\

    \theutterance \stepcounter{utterance}  
    & & & \multicolumn{2}{p{0.3\linewidth}}{
        \cellcolor[rgb]{0.9,0.9,0.9}{
            \makecell[{{p{\linewidth}}}]{
                \texttt{\tiny{[GM$|$GM]}}
                \texttt{no} \\
            }
        }
    }
    & & \\ \\

    \theutterance \stepcounter{utterance}  
    & & & \multicolumn{2}{p{0.3\linewidth}}{
        \cellcolor[rgb]{0.9,0.9,0.9}{
            \makecell[{{p{\linewidth}}}]{
                \texttt{\tiny{[GM$|$GM]}}
                \texttt{game\_result = WIN} \\
            }
        }
    }
    & & \\ \\

\end{supertabular}
}

\end{document}
